% English Version
\documentclass[12pt,a4paper]{article}
\usepackage[utf8]{inputenc}
\usepackage[T1]{fontenc}
\usepackage{lmodern}
\usepackage{amsmath,amssymb}
\usepackage{graphicx}
\usepackage{xcolor}
\usepackage{hyperref}
\usepackage{geometry}
\geometry{a4paper, left=3cm, right=3cm, top=3cm, bottom=3cm}
\usepackage{setspace}
\onehalfspacing
\usepackage{parskip}
\usepackage[english]{babel}
\usepackage{csquotes}
\usepackage{microtype}
\usepackage{booktabs}
\usepackage{longtable}
\usepackage{array}
\usepackage{listings}
\usepackage{caption}
\usepackage{subcaption}
\usepackage{float}
\usepackage{url}
\usepackage{natbib}
\usepackage{titling}

\lstset{
  language=Python,
  basicstyle=\ttfamily\small,
  keywordstyle=\color{blue},
  commentstyle=\color{green!40!black},
  stringstyle=\color{red},
  showstringspaces=false,
  numbers=left,
  numberstyle=\tiny,
  numbersep=5pt,
  breaklines=true,
  frame=single,
  backgroundcolor=\color{gray!5},
  tabsize=2,
  captionpos=b
}

\title{\Huge\textbf{ARS as a Meta-Methodology} \\[2mm]
       \LARGE The Conditions for Explanatory Models \\[2mm]
       \LARGE in the Age of Generative AI}

\author{
  \large
  \begin{tabular}{c}
    Paul Koop
  \end{tabular}
}

\date{\large 1994--2026}

\begin{document}

\maketitle

\begin{abstract}
This paper examines the Algorithmic Recursive Sequence Analysis (ARS) not merely 
as a method but as a \textit{meta-methodology}—a framework that specifies the 
conditions under which a model counts as \textit{explanatory} rather than merely 
\textit{descriptive} or \textit{simulative}. Drawing on the ARS's core principles—the 
primacy of interpretation, the separation of structure and statistics, controlled 
falsification, and XAI validation—I argue that these principles constitute necessary 
criteria for explanatory models in any discipline that deals with sequential social 
processes. The paper systematically relates ARS to five contemporary research 
programs: (1) Formal Verification and Model Checking, (2) Interpretable Machine 
Learning and Rule Extraction, (3) Grounded Theory, (4) Causal Inference, and 
(5) Process Mining. For each, I demonstrate what ARS can learn from these 
approaches and, crucially, what these approaches can learn from ARS. The paper 
concludes that ARS provides a transdisciplinary benchmark for distinguishing 
genuine explanation from statistical or structural description.
\end{abstract}

\newpage
\tableofcontents
\newpage

\section{Introduction: From Method to Meta-Methodology}

The Algorithmic Recursive Sequence Analysis (ARS), in its various versions 
(2.0–4.0), has been presented primarily as a method for the formal analysis 
of sequential interactions. It transforms interpretively obtained categories 
into formal models: probabilistic context-free grammars (PCFG), Petri nets, 
Bayesian networks, and deterministic finite automata (DFA). This paper takes 
a different perspective. It argues that ARS is not just a method but a 
\textit{meta-methodology}—a framework that specifies the \textit{conditions} 
under which a model can be considered explanatory.

This shift in perspective is motivated by a recurring confusion in contemporary 
AI and data science: the conflation of \textbf{simulation} with \textbf{explanation}. 
A large language model (LLM) trained on sales conversations can simulate 
plausible dialogues with high fidelity. Yet, as the ARS notebooks have 
demonstrated, it cannot explain \textit{why} a particular sequence of speech 
acts is well-formed or what rules constitute its generation. The LLM provides 
a statistical shadow of the process; the ARS provides its structural skeleton.

The meta-methodological claim of this paper is that the ARS principles 
constitute \textbf{necessary conditions} for explanatory models in any 
discipline dealing with sequential social processes. These conditions are:

\begin{enumerate}
    \item \textbf{Interpretive Grounding}: Symbols must be tied to documented 
    qualitative interpretations, not merely to statistical correlations.
    
    \item \textbf{Structural Decidability}: The well-formedness of sequences 
    must be formally decidable (e.g., by a DFA), independent of empirical 
    frequencies.
    
    \item \textbf{Generative Transparency}: The model must be able to generate 
    sequences in a way that every step is traceable to explicit rules.
    
    \item \textbf{Falsifiability}: Interpretations must be subject to 
    controlled falsification; counterexamples must be able to refute rules.
    
    \item \textbf{XAI Validation}: The model must satisfy the NIST XAI 
    criteria—meaningfulness, accuracy, and knowledge limits.
\end{enumerate}

To substantiate this claim, I systematically relate ARS to five contemporary 
research programs that share overlapping concerns: formal verification, 
interpretable machine learning, grounded theory, causal inference, and process 
mining. For each, I ask two reciprocal questions:

\begin{enumerate}
    \item What can ARS learn from this approach? (Technical or conceptual 
    enhancements)
    
    \item What can this approach learn from ARS? (Methodological safeguards, 
    criteria for explanation)
\end{enumerate}

\section{ARS as a Meta-Methodology: The Core Principles}

Before examining the five approaches, it is necessary to state the ARS 
principles that serve as the meta-methodological benchmark.

\subsection{Principle 1: Interpretive Grounding}

In ARS, every terminal symbol (KBG, VBBd, KAA, etc.) is the product of a 
documented qualitative interpretation. The coding process is not a black box; 
it is recorded, justified, and subject to intersubjective validation. This 
principle excludes purely data-driven category formation (e.g., clustering 
embeddings) from counting as explanation.

\subsection{Principle 2: Separation of Structure and Statistics}

As developed in \texttt{ARS\_XAI\_Aut2\_Ger.tex}, ARS maintains a strict 
separation between structural rules (which are deterministic and context-free) 
and statistical regularities (which are empirical and contingent). A structural 
rule holds or does not hold—independent of how often it is violated. This 
separation is absent in purely statistical models.

\subsection{Principle 3: Generative Transparency}

The induced grammar must be able to generate sequences in a traceable manner. 
The transducer in Lisp, the parser in Pascal, and the inductor in Scheme each 
provide a different window into this transparency. A model that cannot generate 
exemplars from its own rules is not explanatory.

\subsection{Principle 4: Controlled Falsification}

Interpretations are not asserted once and for all. They are produced as 
readings and then falsified by subsequent sequence positions (following 
Oevermann's sequential analysis). This creates a spiral of interpretation and 
refutation that mirrors Popperian falsificationism adapted to qualitative 
material.

\subsection{Principle 5: XAI Validation}

The three NIST XAI criteria—meaningfulness, accuracy, knowledge limits—are 
not optional add-ons but constitutive elements of explanatory models. 
Meaningfulness requires semantic interpretability; accuracy requires 
correspondence with the material; knowledge limits require explicit 
documentation of the model's boundaries.

\section{Five Approaches in Dialogue with ARS}

\subsection{Formal Verification and Model Checking}

\subsubsection{What Formal Verification Is}

Formal verification, particularly model checking, is a method from theoretical 
computer science that systematically checks whether a formal model (e.g., a 
finite automaton, a Petri net, or a Bayesian network) satisfies specified 
properties. These properties include \textit{safety} ("something bad never 
happens") and \textit{liveness} ("something good eventually happens").

Model checking is exhaustive: it explores the entire state space of the model.
Unlike statistical testing, which provides probabilistic guarantees, model 
checking provides \textit{deterministic} guarantees about the model's behavior.

\subsubsection{What ARS Can Learn from Formal Verification}

\begin{itemize}
    \item \textbf{Property Specification Languages}: ARS could adopt temporal 
    logics (LTL, CTL) to specify what properties the induced grammar should 
    satisfy. For example: $\square (KBG \rightarrow \lozenge VBG)$ ("Always, 
    if a customer greets, eventually the seller greets back").
    
    \item \textbf{Counterexample Generation}: When a property fails, model 
    checkers produce a counterexample trace. This could serve as a powerful 
    falsification tool for ARS interpretations.
    
    \item \textbf{State Space Explosion Awareness}: ARS modelers should be 
    aware that hierarchical grammars can lead to large state spaces. Formal 
    verification offers abstraction techniques to manage this.
\end{itemize}

\subsubsection{What Formal Verification Can Learn from ARS}

\begin{itemize}
    \item \textbf{Interpretive Grounding of States}: In standard model checking, 
    states are abstract symbols. ARS insists that each state must be 
    interpretively grounded. This could lead to a new subfield: \textit{interpretive 
    model checking}, where properties are checked \textit{and} the meaning of 
    states is documented.
    
    \item \textbf{The Separation of Structure and Statistics}: Model checking 
    typically assumes a deterministic model. ARS shows how to separate 
    structural rules (which can be verified) from statistical variations 
    (which cannot). This could enrich probabilistic model checking.
    
    \item \textbf{XAI Criteria for Verification}: Model checking results are 
    often opaque to non-experts. ARS's XAI criteria could guide the development 
    of more understandable verification outputs.
\end{itemize}

\subsection{Interpretable Machine Learning and Rule Extraction}

\subsubsection{What IML and Rule Extraction Are}

Interpretable Machine Learning (IML) aims to make black-box models (neural 
networks, gradient boosting, random forests) understandable to humans. 
\textit{Rule extraction} is a specific IML technique that attempts to describe 
the learned function approximately through a set of if-then rules (e.g., using 
RIPPER, CART, or analyzing activation patterns).

Unlike ARS, which induces rules directly from data, rule extraction typically 
works \textit{post-hoc}: the model is already trained, and rules are extracted 
as an explanation.

\subsubsection{What ARS Can Learn from IML and Rule Extraction}

\begin{itemize}
    \item \textbf{Scalable Rule Induction}: ARS currently induces rules from 
    small corpora (n = 8). IML offers techniques for extracting rules from 
    large datasets, albeit with less interpretive control.
    
    \item \textbf{Quantitative Rule Evaluation}: IML provides metrics for 
    evaluating extracted rules (coverage, fidelity, stability). ARS could 
    adopt these to assess how well a grammar generalizes.
    
    \item \textbf{Hybrid Rule Sets}: Some IML methods combine global and local 
    rules. ARS could explore hybrid grammars that have both a core structural 
    grammar and local statistical variations.
\end{itemize}

\subsubsection{What IML and Rule Extraction Can Learn from ARS}

\begin{itemize}
    \item \textbf{Explication vs. Approximation}: IML's rule extraction is 
    almost always approximate. ARS insists on \textit{exact} rules for the 
    given corpus. This raises a fundamental question: Is approximation ever 
    acceptable for explanation? ARS suggests a clear answer: approximation 
    is acceptable only if the approximation error is documented and the 
    structural core is exact.
    
    \item \textbf{Interpretive Validation of Rules}: IML extracts rules that 
    statistically fit the data. ARS adds a layer of \textit{interpretive 
    validation}: rules must also make sense to human interpreters. This could 
    prevent the extraction of statistically correct but semantically meaningless 
    rules.
    
    \item \textbf{Falsifiability as a Criterion}: IML rarely discusses how 
    extracted rules could be falsified. ARS makes falsifiability a central 
    criterion. IML could adopt this: a rule set is not explanatory if no 
    conceivable counterexample could refute it.
\end{itemize}

\subsection{Grounded Theory}

\subsubsection{What Grounded Theory Is}

Grounded Theory (GT), developed by Glaser and Strauss, is a classic methodology 
in qualitative social research for developing theories from data. It involves 
procedures such as open coding, axial coding, and selective coding. The goal 
is to generate middle-range theories that are "grounded" in empirical material.

In recent years, there have been attempts to formalize parts of GT or to 
support it computationally (e.g., through natural language processing or 
topic modeling). However, GT remains largely informal in its final output.

\subsubsection{What ARS Can Learn from Grounded Theory}

\begin{itemize}
    \item \textbf{Systematic Coding Procedures}: GT offers a rich vocabulary 
    and set of procedures for coding that could enrich ARS's interpretive 
    phase. Concepts like "axial coding" (relating categories to subcategories) 
    are similar to ARS's hierarchical compression but more fine-grained.
    
    \item \textbf{Theoretical Sampling}: GT's principle of theoretical 
    sampling—selecting new cases based on emerging theoretical insights—could 
    guide ARS's case selection in larger studies.
    
    \item \textbf{Constant Comparison}: GT's method of constant comparison 
    (comparing each new case with already developed categories) is already 
    implicit in ARS's systematic case comparison (Phase 4) but could be 
    made more explicit.
\end{itemize}

\subsubsection{What Grounded Theory Can Learn from ARS}

\begin{itemize}
    \item \textbf{From Theory to Generative Model}: GT typically stops at the 
    level of narrative theory or category systems. ARS goes further: it 
    transforms the theory into a \textit{generative grammar} that can produce 
    new sequences. GT could adopt this to make its theories testable and 
    executable.
    
    \item \textbf{Formal Falsifiability}: GT's validation procedures are 
    primarily qualitative (e.g., member checking, peer debriefing). ARS adds 
    formal falsifiability: the grammar can be wrong in a way that can be 
    demonstrated mechanically (e.g., by a parser rejecting a sequence).
    
    \item \textbf{XAI Criteria for Grounded Theory}: ARS's XAI criteria 
    (meaningfulness, accuracy, knowledge limits) provide a checklist for 
    evaluating grounded theories. A GT theory that cannot specify its 
    knowledge limits is incomplete.
\end{itemize}

\subsection{Causal Inference and Causal Graphical Models}

\subsubsection{What Causal Inference Is}

Causal inference, particularly with causal graphical models (e.g., DAGs, 
DoWhy, CausalNex), goes beyond mere correlation and attempts to identify 
and quantify causal relationships between variables. It uses techniques 
such as the do-calculus, instrumental variables, and counterfactual reasoning.

A key insight of causal inference is that correlation is not causation. 
Directed acyclic graphs (DAGs) are used to represent assumptions about 
causal structure.

\subsubsection{What ARS Can Learn from Causal Inference}

\begin{itemize}
    \item \textbf{Causal Interpretation of Grammars}: ARS grammars describe 
    sequential dependencies. Causal inference could help distinguish whether 
    these dependencies are merely sequential or genuinely causal. For example: 
    Does the question "Anything else?" \textit{cause} an additional purchase, 
    or is it merely correlated with it?
    
    \item \textbf{Counterfactual Reasoning}: Causal inference excels at 
    answering counterfactual questions ("What would have happened if the 
    seller had not asked?"). ARS could adopt counterfactual simulation 
    (already present in Phase 3 of CGTI) as a standard validation tool.
    
    \item \textbf{Instrumental Variables for Sequential Data}: ARS deals with 
    sequential data where confounding is common. Instrumental variable 
    techniques could help identify causal effects even in observational 
    sequential data.
\end{itemize}

\subsubsection{What Causal Inference Can Learn from ARS}

\begin{itemize}
    \item \textbf{Interpretive Grounding of Causal Graphs}: In standard causal 
    inference, the DAG is often assumed or learned from data without interpretive 
    documentation. ARS insists that every node and edge must be interpretively 
    grounded. This could lead to \textit{interpretive causal inference} as a 
    new subfield.
    
    \item \textbf{Sequential Grammars as Causal Structures}: ARS grammars are 
    a form of causal structure over sequences. Causal inference typically 
    deals with static or time-series data, not with grammatical sequences. 
    ARS could inspire a new class of \textit{grammatical causal models}.
    
    \item \textbf{XAI for Causal Models}: Causal models are often presented 
    as DAGs with probabilities, which are not self-explanatory. ARS's XAI 
    criteria could guide the documentation of causal models, making them 
    more accessible to domain experts.
\end{itemize}

\subsection{Process Mining}

\subsubsection{What Process Mining Is}

Process mining is a research field at the intersection of data mining, machine 
learning, and process modeling. It aims to discover, conform, and enhance 
process models (often in the form of Petri nets, BPMN diagrams, or directly 
follows graphs) from event logs—sequential recordings of process steps, e.g., 
in workflow management systems or ERP systems.

Process mining typically works with large, anonymized, and weakly annotated 
logs. It discovers the "average process" or the "most common paths." It does 
not aim for a complete reconstruction of a single case's constitutive rules.

\subsubsection{What ARS Can Learn from Process Mining}

\begin{itemize}
    \item \textbf{Scalable Discovery Algorithms}: Process mining offers 
    sophisticated algorithms (e.g., Alpha miner, Heuristics miner, Inductive 
    miner) for discovering Petri nets from large logs. ARS could adopt or 
    adapt these for larger corpora while preserving interpretive control.
    
    \item \textbf{Conformance Checking}: Process mining includes techniques 
    for checking whether an event log conforms to a given model. This could 
    be used to validate ARS grammars against new data.
    
    \item \textbf{Performance Analysis}: Process mining adds performance 
    dimensions (time, cost, frequency). ARS could be extended to incorporate 
    temporal and resource dimensions more systematically.
\end{itemize}

\subsubsection{What Process Mining Can Learn from ARS}

\begin{itemize}
    \item \textbf{Interpretive Discovery for Small Logs}: Process mining 
    typically requires large logs to produce reliable models. ARS shows how 
    to discover models from a single case (n = 1) through interpretive 
    depth. This could be valuable for process mining in domains where data 
    is scarce (e.g., medical procedures, legal cases).
    
    \item \textbf{Documentation of Discovery Decisions}: Process mining 
    algorithms make many decisions (e.g., which paths to include, how to 
    handle noise). These decisions are rarely documented in an interpretively 
    accessible way. ARS's reflexive documentation could serve as a model.
    
    \item \textbf{Separation of Structure and Statistics}: Process mining 
    often produces models that mix structural rules with statistical noise. 
    ARS's strict separation could improve the quality of discovered models 
    by clearly distinguishing what is structurally necessary from what is 
    merely empirically frequent.
    
    \item \textbf{XAI for Process Mining}: The models produced by process 
    mining (e.g., spaghetti-like Petri nets) are often hard to understand. 
    ARS's XAI criteria could guide the development of more explainable 
    process mining outputs.
\end{itemize}

\section{Toward a Transdisciplinary Benchmark for Explanation}

The five comparisons above reveal a common pattern. Each contemporary 
approach has technical strengths that could enhance ARS: scalability 
(IML, Process Mining), formal rigor (Verification, Causal Inference), 
and systematic coding procedures (Grounded Theory). Conversely, each 
approach lacks some of the meta-methodological safeguards that ARS 
provides: interpretive grounding, structural decidability, generative 
transparency, controlled falsification, and XAI validation.

This symmetry suggests that ARS is not merely one method among others 
but a \textbf{transdisciplinary benchmark} for what counts as an 
explanatory model.

\begin{table}[H]
\centering
\caption{ARS as a Transdisciplinary Benchmark}
\label{tab:benchmark}
\begin{tabular}{@{} p{3cm} p{4cm} p{6cm} @{}}
\toprule
\textbf{Criterion} & \textbf{Question} & \textbf{Absence indicates...} \\
\midrule
Interpretive Grounding & Are the symbols meaningfully documented? & Statistical correlation without understanding \\
Structural Decidability & Is well-formedness formally decidable? & Probabilistic guessing rather than rule-following \\
Generative Transparency & Can the model generate exemplars traceably? & Simulation without explanation \\
Controlled Falsification & Can counterexamples refute rules? & Unfalsifiable post-hoc storytelling \\
XAI Validation & Are meaningfulness, accuracy, and knowledge limits documented? & Technical sophistication without epistemic accountability \\
\bottomrule
\end{tabular}
\end{table}

\subsection{The Explanation vs. Simulation Distinction Revisited}

The core distinction that emerges from this analysis is between 
\textbf{explanation} and \textbf{simulation}. A model \textit{simulates} 
if it reproduces the statistical properties of the data. A model 
\textit{explains} if it specifies the constitutive rules that generate 
the phenomenon.

\begin{itemize}
    \item \textbf{Simulation} is sufficient for prediction. An LLM that 
    accurately predicts the next token is a good simulator.
    \item \textbf{Explanation} is necessary for understanding, intervention, 
    and normative evaluation. An ARS grammar that specifies the rules of 
    a sales conversation is an explanation.
\end{itemize}

The five criteria above are the conditions under which a model qualifies 
as an explanation rather than merely a simulation.

\subsection{The Role of the Human Interpreter}

A recurring theme across all five comparisons is the role of the human 
interpreter. In ARS, the human is constitutive: interpretation is a human 
act that cannot be fully automated. In the other approaches, the human 
is often external—designing algorithms, providing training data, evaluating 
outputs.

This is not a weakness of ARS but a strength. ARS makes explicit what is 
often implicit: that explanation is a human practice, not a property of a 
model considered in isolation. A model is explanatory \textit{for someone} 
who can understand it, use it, and be held accountable for it.

\section{Conclusion and Outlook}

This paper has argued that the Algorithmic Recursive Sequence Analysis 
(ARS) is not merely a method but a meta-methodology—a framework that 
specifies the conditions under which a model counts as explanatory. 
Five core principles were identified: interpretive grounding, structural 
decidability, generative transparency, controlled falsification, and XAI 
validation.

The paper then systematically related ARS to five contemporary research 
programs: formal verification, interpretable machine learning, grounded 
theory, causal inference, and process mining. For each, I demonstrated 
what ARS can learn from these approaches (technical enhancements) and, 
crucially, what these approaches can learn from ARS (methodological 
safeguards, criteria for explanation).

The meta-methodological claim is that the ARS principles constitute 
necessary conditions for explanatory models in any discipline dealing 
with sequential social processes. They provide a transdisciplinary 
benchmark for distinguishing genuine explanation from statistical or 
structural description.

For future research, three directions are particularly promising:

\begin{enumerate}
    \item \textbf{Implementation of hybrid systems}: Integrate ARS 
    grammars with model checkers, rule extractors, or process mining 
    algorithms while preserving the meta-methodological safeguards.
    
    \item \textbf{Empirical testing of the benchmark}: Apply the five 
    criteria to existing models in different disciplines and assess 
    whether they predict perceived explanatory quality.
    
    \item \textbf{Extension to non-sequential domains}: While ARS was 
    developed for sequential data, the meta-methodological principles 
    may generalize to other types of models (e.g., classification, 
    clustering, regression).
\end{enumerate}

In conclusion: The question is not whether a model fits the data. 
Statistical fit is necessary but not sufficient. The question is whether 
the model meets the meta-methodological criteria that make explanation 
possible. ARS provides a concrete, operationalized answer.

\newpage
\begin{thebibliography}{99}

\bibitem[Baier \& Katoen(2008)]{baier2008principles}
Baier, C., \& Katoen, J.-P. (2008). \textit{Principles of Model Checking}. 
MIT Press.

\bibitem[Glaser \& Strauss(1967)]{glaser1967discovery}
Glaser, B. G., \& Strauss, A. L. (1967). \textit{The Discovery of Grounded 
Theory: Strategies for Qualitative Research}. Aldine.

\bibitem[Koop(1992)]{koop1992parser}
Koop, P. (1992). \textit{Demo-Parser Chart-Parser Version 1.0}. Pascal source code.

\bibitem[Koop(1994)]{koop1994scheme}
Koop, P. (1994). \textit{Grammatikinduktion empirisch gesicherter 
Verkaufsgespräche}. Scheme source code.

\bibitem[Koop(1994)]{koop1994lisp}
Koop, P. (1994). \textit{Sequenzanalyse empirisch gesicherter 
Verkaufsgespräche}. Lisp source code.

\bibitem[Koop(2023)]{koop2023notebook}
Koop, P. (2023). \textit{Qualitative Sozialforschung und Große Sprachmodelle}. 
Jupyter Notebook.

\bibitem[Koop(2024)]{koop2024ars}
Koop, P. (2024/2026). \textit{Between Interpretation and Computation: 
Algorithmic Recursive Sequence Analysis as a Bridge between Qualitative 
Hermeneutics and Formal Modeling}. the-last-freedom.org.

\bibitem[Koop(2026)]{koop2026causal}
Koop, P. (2026). \textit{Social Structures and Processes: Causal Inference 
with Probabilistic Context-Free Grammars and Bayesian Networks}. 
the-last-freedom.org.

\bibitem[Molnar(2022)]{molnar2022interpretable}
Molnar, C. (2022). \textit{Interpretable Machine Learning: A Guide for 
Making Black Box Models Explainable}. leanpub.com.

\bibitem[Oevermann et al.(1979)]{oevermann1979methodology}
Oevermann, U., Allert, T., Konau, E., \& Krambeck, J. (1979). The methodology 
of objective hermeneutics. In H.-G. Soeffner (Ed.), \textit{Interpretative 
Procedures in the Social and Text Sciences} (pp. 352-434). Metzler.

\bibitem[Ortigossa et al.(2024)]{ortigossa2024xai}
Ortigossa, E. S., Gonçalves, T., \& Nonato, L. G. (2024). Explainable 
Artificial Intelligence (XAI)—From Theory to Methods and Applications. 
\textit{IEEE Access}, 12, 80799-80846.

\bibitem[Pearl(2009)]{pearl2009causality}
Pearl, J. (2009). \textit{Causality: Models, Reasoning, and Inference} 
(2nd ed.). Cambridge University Press.

\bibitem[van der Aalst(2016)]{vanderaalst2016process}
van der Aalst, W. M. P. (2016). \textit{Process Mining: Data Science in 
Action} (2nd ed.). Springer.

\end{thebibliography}

\end{document}